\documentclass[10pt,a4paper]{amsart}
\usepackage[margin=19mm]{geometry}
\usepackage{amsmath,amssymb,mathtools}
\usepackage[scr=rsfs,cal=euler]{mathalfa}
\usepackage{xfrac}
\usepackage[unicode,hidelinks,bookmarks=false]{hyperref}
\newcommand{\ie}{i.e.\ }
\newcommand{\cf}{cf.\ }
\newcommand{\resp}{respectively\ }
\newcommand{\Subsection}{\subsection}
\providecommand{\nobreakdash}{\nobreak\hbox{-}}
\setlength{\emergencystretch}{2em}
\multlinegap0pt
\usepackage{mathtools}
\usepackage{amssymb}
\newtheorem{lemma}{Lemma}
\title{An optimal Poincaré inequality\\for the complex Ginibre log-gas}
\author{Djalil Chafaï}
\date{Source: arXiv:2608.19358v2 (CC BY 4.0)}
\begin{document}
\maketitle

\noindent\emph{Source and licence.} An optimal Poincaré inequality\\for the complex Ginibre log-gas. Version arXiv:2608.19358v2 (CC BY 4.0), distributed by arXiv under the Creative Commons Attribution 4.0 International licence, http://creativecommons.org/licenses/by/4.0/, as recorded in the arXiv licence metadata for this exact version and shown on the abstract page https://arxiv.org/abs/2608.19358v2. Full licence notice: \texttt{licenses/arxiv-2608.19358-CC-BY-4.0.txt}.\par\smallskip

\noindent\textit{Editorial scope.} Counted unit: the article's lemma le:projection-geometry (source0.tex lines 1634--1667). The article's complete original proof of it is delivered with it (source0.tex lines 1669--1715, one \texttt{proof} environment of 47 lines). No mathematics is written, repaired or re-derived: the statement and the proof are the article's own lines, in the article's own notation, and no retained line is substituted or re-flowed.  Retained uncounted as the source-local context the proof needs: lines 1533--1550.  Nothing was omitted from any retained span, so the package carries no omission record.  No retained line contains a citation, so the wrapper carries no bibliography and no bibliography entry.  No retained line contains a figure environment, an image-inclusion command or drawing code, so no image is delivered and the package declares no asset dependency.  Editorial formatting only, in addition: an AMS \texttt{amsart} preamble that restates the article's own declarations and macros used by the retained lines, namely the environments \texttt{lemma} on separate counters, together with the standard packages those lines need.  The retained environments print in this document as Lemma 1 in order of appearance, whereas the article prints the same items with the numbering of its own full document.  The complete native source of the article as distributed in the arXiv e-print for this exact version is bundled unchanged as \texttt{sources/208080/source0.tex}.
\begin{lemma}[Divisibility]\label{le:divisibility}
  Let $\mathcal{H}_n^{\mathrm{asym}}$ be the subspace of $\mathcal{H}_n$ of
  antisymmetric elements of $\mathcal{H}_n$. Then
  \begin{equation}
    \mathcal{H}_n^{\mathrm{asym}}=V_n\mathcal{H}_n^{\mathrm{div}}
    \quad\text{where}\quad
    \mathcal{H}_n^{\mathrm{div}} = \{h\in
    L^2(\mu_n):V_nh\in\mathcal{H}_n^{\mathrm{asym}}\}.
  \end{equation}
  Equivalently, for every antisymmetric holomorphic $g\in L^2(\gamma_n)$, we have
  $g=V_nf$ for a symmetric holomorphic $f\in L^2(\mu_n)$. Moreover, for every
  symmetric $f\in L^2(\mu_n)$,
  \begin{equation}%\label{eq:distance-transfer}
    \operatorname{dist}_{L^2(\gamma_n)}(V_nf,\mathcal{H}_n)^2
    =
    c_n'\operatorname{dist}_{L^2(\mu_n)}(f,\mathcal{H}_n^{\mathrm{div}})^2.
  \end{equation}
\end{lemma}

\begin{lemma}[Holomorphic--antiholomorphic projection geometry]
  \label{le:projection-geometry}
  The following properties hold.
  \begin{enumerate}
  \item
    $\mathcal{H}_n^{\mathrm{div},0}$ and
    $\overline{\mathcal{H}_n^{\mathrm{div},0}}$ are orthogonal.
  \item For every symmetric $f\in L^2(\mu_n)$,
    \begin{equation}\label{eq:two-projection-bound}
      \|Pf\|_{L^2(\mu_n)}^2+\|Qf\|_{L^2(\mu_n)}^2
      \leq
      \|f\|_{L^2(\mu_n)}^2+\|P_0f\|_{L^2(\mu_n)}^2.
    \end{equation}
  \item If $f\in L^2(\mu_n)$ is real-valued, symmetric, and centered, then
    $h=Pf$, its pointwise complex conjugate $\overline h$, and
    $r=f-h-\overline h$, are three pairwise orthogonal functions in $L^2(\mu_n)$, and
    \begin{align}
      \|f\|_{L^2(\mu_n)}^2
      &=2\|h\|_{L^2(\mu_n)}^2+\|r\|_{L^2(\mu_n)}^2,
      \label{eq:projection-geometry-decomposition}\\
      \operatorname{dist}_{L^2(\mu_n)}
      \bigl(f,\mathcal{H}_n^{\mathrm{div}}\bigr)^2
      &=\frac{1}{2}\|f\|_{L^2(\mu_n)}^2
        +\frac{1}{2}\|r\|_{L^2(\mu_n)}^2.
      \label{eq:exact-half-distance}
    \end{align}
    In particular, we have the half-distance estimate 
    \begin{equation}\label{eq:half-distance-paper}
      \operatorname{dist}_{L^2(\mu_n)}
      \bigl(f,\mathcal{H}_n^{\mathrm{div}}\bigr)^2
      \geq\frac{1}{2}\|f\|_{L^2(\mu_n)}^2.
    \end{equation}
  \end{enumerate}
\end{lemma}

\begin{proof}
  For $\theta\in\mathbb{R}$, let
  $T_\theta f(z)=f(\mathrm{e}^{\mathrm{i}\theta}z)$. Since $\mu_n$ is
  invariant under global rotations, ${(T_\theta)}_{\theta\in\mathbb{R}}$ is a
  group of unitary transformations of $L^2(\mu_n)$. Let
  $h\in\mathcal{H}_n^{\mathrm{div}}$, set $G=V_nh$, and write
  $G=\sum_{m\geq0}G_m$ for the homogeneous expansion of $G$ in
  $\mathcal{H}_n$. Each $G_m$ is antisymmetric. By the divisibility Lemma
  \ref{le:divisibility}, each nonzero $G_m$ is divisible by $V_n$. Since $V_n$
  is homogeneous of degree $d_n=n(n-1)/2$, this gives $G_m=V_nh_{m-d_n}$ for a
  symmetric homogeneous holomorphic polynomial $h_{m-d_n}$ of degree $m-d_n$,
  in particular, $G_m=0$ for $m<d_n$. The Vandermonde isometry yields
  \[
    h=\sum_{k\geq0}h_k
    \quad\text{in }L^2(\mu_n),
  \]
  where $h_k$ has rotation weight $k$. Thus $\mathcal{H}_n^{\mathrm{div},0}$
  is the orthogonal sum of the strictly positive rotation-weight spaces,
  whereas $\overline{\mathcal{H}_n^{\mathrm{div},0}}$ is the orthogonal sum of
  the strictly negative rotation-weight spaces. This proves the first
  assertion. Let $P_+$ and $P_-$ be the orthogonal projections onto these two
  centered subspaces. Then $P=P_0+P_+$ and $Q=P_0+P_-$. To get
  \eqref{eq:two-projection-bound}, we observe that the ranges of $P_0$, $P_+$,
  and $P_-$ are pairwise orthogonal, hence the Bessel inequality gives
  \begin{align*}
    \|Pf\|_2^2+\|Qf\|_2^2
    &=2\|P_0f\|_2^2+\|P_+f\|_2^2+\|P_-f\|_2^2\\
    &\leq2\|P_0f\|_2^2+\|f-P_0f\|_2^2
      =\|f\|_2^2+\|P_0f\|_2^2.
  \end{align*}

  Let $C$ be complex conjugation. It is an antiunitary involution and $Q=CPC$.
  If $f$ is real-valued and centered, then $Qf=C(Pf)=\overline h$ and
  $P_0f=0$. Hence $h$, $\overline h$, and $r=f-h-\overline h$ are the
  orthogonal components of $f$ in the ranges of $P_+$, $P_-$, and their common
  orthogonal complement. This proves
  \eqref{eq:projection-geometry-decomposition}. Finally,
  \eqref{eq:exact-half-distance} follows from
  \eqref{eq:projection-geometry-decomposition} since
  \[
    \operatorname{dist}_{L^2(\mu_n)}
    \bigl(f,\mathcal{H}_n^{\mathrm{div}}\bigr)^2
    =\|f-h\|_2^2
    =\|\overline h+r\|_2^2
    =\|h\|_2^2+\|r\|_2^2.
  \]
\end{proof}

\end{document}
